\section{Overview}
\label{sec:overview}

codetopo is a structural code intelligence MCP server for AI agents. It indexes
source trees with tree-sitter~\cite{brunsfeld2018treesitter}, builds a persistent
symbol graph in SQLite, and exposes that graph to language models and developer
<<<<<<< HEAD
tools through the Model Context Protocol~\cite{anthropic2024mcp}.\footnote{MCP was publicly announced at Anthropic's developer day in November 2024. The protocol version string \texttt{2024-11-05} reflects this origin date; MCP versions are identified by calendar date.}
=======
tools through the Model Context Protocol~\cite{anthropic2024mcp}.
>>>>>>> origin/main

\subsection{The central claim}

Language models that navigate codebases today overwhelmingly rely on text
search: grep-style substring matching or embedding-vector approximate retrieval.
Both approaches treat source code as a bag of tokens. codetopo takes a
different stance: source code has \emph{structure}, and that structure---the
symbol graph---is the right primary index.

A symbol graph records what symbols exist (\emph{nodes}: functions, classes,
variables, types), how they relate (\emph{edges}: calls, imports, inheritance,
field access), and where they live (\emph{location}: file, line, column). Given
a symbol graph an LLM can answer questions that text search cannot:
``what calls \texttt{parse\_expr}?'', ``what does \texttt{Widget} inherit from?'',
``what symbols does \texttt{server.cpp} export?'' --- in milliseconds, without
reading source files. Figure~\ref{fig:symbol-graph} shows the scale of graph
that codetopo treats as its fundamental unit of navigation.

\begin{figure}[htbp]
  \centering
  \tritonnext{width=0.78\textwidth}
\begin{triton}
nodegraph
  directed
  title Tiny codetopo symbol graph
  node IDX : index_file()
  node PAR : parse_file()
  node INS : insert_batch()
  node SRV : mcp_server()
  IDX -> PAR : calls
  IDX -> INS : writes
  SRV -> IDX : imports
\end{triton}
  \caption{A small codetopo symbol graph with symbol nodes and labeled call,
           import, and write edges.}
  \label{fig:symbol-graph}
\end{figure}

\subsection{The problem it solves}

Current LLM-assisted coding tools face a fundamental tension: context windows
are finite, but codebases are large. The standard response is to retrieve
snippets by embedding similarity. This works for local, semantic queries but
fails for structural ones. A call-graph traversal expressed as a text search
requires the model to read every file that might contain a call; with codetopo
the traversal is a single SQL join.

codetopo shifts the \emph{navigation burden} from the model to the index.
The model asks a precise structural question (``give me the callers of
\texttt{foo} up to depth~3'') and receives a compact, structured answer. It
never needs to load intermediate files. This design is motivated by empirical
findings on how developers navigate large codebases~\cite{robillard2004codenav}.

\subsection{Design properties}

\begin{description}
  \item[Structural-first.] The primary index is a symbol graph, not a text
        corpus. Full-text search is provided as a complement via an FTS5
        trigram index~\cite{hipp2014fts5}, not as a substitute.
  \item[Zero-dependency runtime.] The MCP server is a single compiled binary;
        the index is a single SQLite file. No language server, no daemon, no
        network service is required.
  \item[Polyglot.] Thirteen languages are supported out of the box: C, C++,
        C\#, TypeScript, JavaScript, Python, Rust, Java, Go, Bash, SQL, and
<<<<<<< HEAD
        YAML.\footnote{Lua is the thirteenth language; it is tracked as
        experimental and omitted from this list. Simon's architecture notes
        enumerate twelve production-ready languages.} Any language with a
        tree-sitter grammar can be added with a thin extractor.
=======
        YAML. Any language with a tree-sitter grammar can be added with a
        thin extractor.
>>>>>>> origin/main
  \item[Multi-root.] A single server instance manages multiple workspace
        roots simultaneously (\S\ref{sec:workspace}).
  \item[Incremental.] Files are re-indexed only when their content changes
        (mtime $+$ xxHash check). A filesystem watcher triggers live
        re-indexing on create, modify, and delete events.
  \item[MCP-native.] codetopo speaks the Model Context Protocol
        (\S\ref{sec:mcp-protocol}) natively: JSON-RPC~2.0 over stdio, the
<<<<<<< HEAD
        same transport as         the Language Server Protocol~\cite{microsoft2016lsp}\footnote{LSP was
        first announced by Microsoft and Red Hat on June~27, 2016 ---
        coincidentally the same calendar date as this document's writing.}
=======
        same transport as the Language Server Protocol~\cite{microsoft2016lsp}
>>>>>>> origin/main
        but with a tool-call interface designed for autonomous agents.
\end{description}

\subsection{What codetopo is not}

codetopo is a code intelligence \emph{index and query server}. It does
\emph{not} compile code, does \emph{not} run tests, does \emph{not} generate
code, and does \emph{not} replace a full language server. It provides the
structural substrate on which tools and agents can build.
